\input{../tex-inputs/latex-leseansicht-vorspann}

\section[Arthur Schnitzler an Rainer Maria Rilke, 17. 1. 1902]{L04345 Arthur Schnitzler an Rainer Maria Rilke, 17. 1. 1902}
\nopagebreak\mylabel{L04345v}
\rehead{ }\normalsize\beginnumbering\briefempfaengerindex{Rilke, Rainer Maria@\textsc{Rilke, Rainer Maria}!zzzSchnitzler, Arthur@\emph{von Arthur Schnitzler}!1902-01-171@{17. Januar 1902}|(be}
\toendnotes[C]{\smallbreak\pagebreak[2]}
\correspDesc{Versand  durch Arthur Schnitzler am 17. Januar 1902 in Wien
\newline{}Erhalt  durch Rainer Maria Rilke im Zeitraum [17. 1. 1902 – 20. 1. 1902?] in Wien}\toendnotes[C]{\smallbreak}
\Standort{DLA, A:Schnitzler, HS.2022.81 XXXX.}
\physDesc{Brief, 1 Blatt, 2 Seiten, 747 Zeichen
\newline{}Handschrift: schwarze Tinte, deutsche Kurrent}
\pstart
           \raggedleft{}{\pb}17. 1. 902\pend
           \vspace{0.5em}
\pstart
           lieber Herr \textsc{Rilke}, ich nehme es
      als einen willko{\geminationm}enen Beweis an
      Freundſchaft, daſs Sie mir{ }ſo raſch
      von den Dingen, die{ }ſich zugetragen haben,
      Mittheilung machen. Inwieweit
      ich Ihnen von Nutzen{ }ſein kann,
      – ob überhaupt, weiſs ich noch nicht;
      aber es wäre wohl denkbar, daſs da
      oder dort ein Wort von mir angehört werden könnte und ſeine Wirkung thut. Wenn Sie nichts da{\pb}gegen haben,{ }ſpreche ich nächſtens mit
              Bahr\pwindex{Bahr, Hermann 19.\,7.\,1863 Linz – 15.\,1.\,1934 München@\textsc{Bahr, Hermann} (19.\,7.\,1863 Linz – 15.\,1.\,1934 München), \emph{Schriftsteller, Kritiker}|pw}; eventuell auch mit Singer\pwindex{Singer, Isidor 16.\,1.\,1857 Budapest – 8.\,12.\,1927 Wien@\textsc{Singer, Isidor} (16.\,1.\,1857 Budapest – 8.\,12.\,1927 Wien), \emph{Journalist, Herausgeber, Soziologe}|pw}
               (Zeit\orgindex{Zeit. Wiener Wochenschrift [Redaktion]@Die Zeit. Wiener Wochenschrift [Redaktion]|pw}). Mit den Tagesblättern\orgindex{Neues Wiener Tagblatt@Neues Wiener Tagblatt|pw}
      habe ich{ }ſo gut wie keine Verb\textcolor{gray}{in}dung. Hoffentlich wird die
              Trennung von Ihrer Familie\pwindex{Sieber, Ruth 12.\,12.\,1901 Westerwede – 27.\,11.\,1972@\textsc{Sieber, Ruth} (12.\,12.\,1901 Westerwede – 27.\,11.\,1972)|pw}\pwindex{Westhoff, Clara 21.\,11.\,1878 Bremen – 9.\,3.\,1954 Fischerhude@\textsc{Westhoff, Clara} (21.\,11.\,1878 Bremen – 9.\,3.\,1954 Fischerhude), \emph{Malerin, Bildhauerin}|pw} nicht
      allzulange dauern – das{ }ſcheint
      mir das weſentlichſte. Ich höre
      jedenfalls{ }ſehr bald von Ihnen, insbeſondre ob Sie, \textsc{resp}. wann Sie
      nach Wien\oindex{Wien@\textbf{Wien}, \emph{Verwaltungsgebiet}|pw} kommen.\pend
           \pstart Herzlichſt der Ihrige,\spacefill\mbox{Arth Schnitzler}\pend{}\selectlanguage{ngerman}\endnumbering\briefempfaengerindex{Rilke, Rainer Maria@\textsc{Rilke, Rainer Maria}!zzzSchnitzler, Arthur@\emph{von Arthur Schnitzler}!1902-01-171@{17. Januar 1902}|)be}\mylabel{L04345h}
\begin{anhang}
\end{anhang}\newcommand{\dateiname}{L04345}\newcommand{\titel}{Arthur Schnitzler an Rainer Maria Rilke, 17. 1. 1902}\newcommand{\editorInnen}{Herausgegeben von Jahnke, SelmaMüller, Martin Anton}\input{../tex-inputs/latex-leseansicht-abspann}
